\subsection{Luồng xử lý verdict, reveal và restart}
\label{sec:engine_flow}
\begin{sloppypar}

\texttt{GameEngine} cung cấp ba thao tác công khai là \texttt{get\_public\_state},
\texttt{submit\_verdict} và \texttt{restart}. Mục này mô tả luồng thực thi của
các thao tác đó, còn điều kiện logic tương ứng với từng kết quả đã được hình
thức hóa ở Mục~\ref{sec:verdict_acceptance} (Bảng~\ref{tab:submit_outcome}).

\paragraph{Khởi tạo}
Hàm khởi tạo nhận vào một \texttt{PuzzleDefinition} cùng một bộ kiểm tra, tính
trước tập định danh ô hợp lệ dưới dạng \texttt{frozenset} để việc kiểm tra một ô
có tồn tại hay không chỉ tốn $O(1)$, rồi gọi \texttt{\_initial\_state()}. Hàm
này tạo một \texttt{GameState} rỗng và áp dụng lần lượt các phán quyết trong
\texttt{initial\_verdicts} thông qua chính hàm \texttt{\_apply\_proof} vốn dùng
cho phán quyết thông thường. Nhờ vậy, nhân vật khởi đầu vừa được đánh dấu là đã
chứng minh, vừa có lời khai được lật sẵn, mà không cần viết thêm mã xử lý riêng
cho trạng thái ban đầu. Lưu ý rằng engine không gọi bộ kiểm tra đối với các
phán quyết ban đầu, vì tính đúng đắn của chúng đã được \texttt{validate\_puzzle}
xác nhận từ trước khi đối chiếu với đáp án.

\paragraph{Nộp phán quyết}
Luồng xử lý của \texttt{submit\_verdict(cell\_id, status)} được trình bày trong
Thuật toán~\ref{alg:submit_verdict}. Hai bước kiểm tra đầu tiên không cần đến
bộ giải nên được thực hiện trước; sau đó engine gọi \texttt{has\_solution} đúng
hai lần trên \emph{cùng một} ảnh chụp công khai, một lần với giả thiết của người
chơi và một lần với giả thiết ngược lại. Kết quả trong mọi trường hợp đều được
đóng gói thành \texttt{SubmitResult(verdict, outcome, public\_state)}, và engine
không ném ngoại lệ cho nước đi sai. Việc ánh xạ kết quả sang mã HTTP được giao cho
API, trong đó hiện chỉ có \texttt{INVALID\_CELL} được chuyển thành HTTP~404, còn
năm kết quả còn lại được trả về với HTTP~200 kèm trường \texttt{outcome}.

\begin{algorithm}[htbp]
  \caption{\texttt{GameEngine.submit\_verdict}}
  \label{alg:submit_verdict}
  \begin{algorithmic}[1]
    \Require ô $i$, trạng thái $s$
    \Ensure \texttt{SubmitResult}
    \State $P \gets \Call{GetPublicState}{}$ \Comment{ảnh chụp trước nước đi}
    \If{$i \notin$ tập ô hợp lệ} \Return $(\langle i,s\rangle, \texttt{INVALID\_CELL}, P)$ \EndIf
    \If{$i \in \texttt{proved\_statuses}$} \Return $(\langle i,s\rangle, \texttt{ALREADY\_PROVED}, P)$ \EndIf
    \State $a \gets \textsc{HasSolution}(P, \langle i, s\rangle)$
    \State $b \gets \textsc{HasSolution}(P, \langle i, \neg s\rangle)$
    \If{$a \land b$} \Return $(\cdot, \texttt{NOT\_PROVABLE}, P)$ \EndIf
    \If{$\neg a \land b$} \Return $(\cdot, \texttt{CONTRADICTED}, P)$ \EndIf
    \If{$\neg a \land \neg b$} \Return $(\cdot, \texttt{INCONSISTENT}, P)$ \EndIf
    \State $\Call{ApplyProof}{\mathit{state}, i, s}$ \Comment{chỉ nhánh $a \land \neg b$ tới đây}
    \State \Return $(\langle i,s\rangle, \texttt{ACCEPTED}, \Call{GetPublicState}{})$
  \end{algorithmic}
\end{algorithm}

Từ thuật toán có thể thấy chỉ kết quả \texttt{ACCEPTED} mới ghi vào trạng thái
ván. Năm kết quả còn lại đều trả về chính ảnh chụp $P$ được dựng ở dòng~1, nên
một nước đi sai không để lại bất kỳ dấu vết nào trong \texttt{GameState}; các ca
kiểm thử \texttt{test\_submit\_*} khẳng định điều này bằng cách so sánh
\texttt{engine.get\_public\_state()} với ảnh chụp trước nước đi. Bên cạnh đó,
\texttt{submit\_verdict} hoàn toàn không đọc \texttt{PuzzleDefinition.solution},
vì vậy một phán quyết dù đúng với đáp án nhưng chưa suy ra được từ thông tin
công khai vẫn nhận kết quả \texttt{NOT\_PROVABLE}. Engine cũng không dùng bộ nhớ
đệm: mỗi lần nộp tốn đúng hai lời gọi \texttt{has\_solution}, và bộ kiểm tra dựng
lại toàn bộ ràng buộc từ ảnh chụp ở mỗi lần gọi. Cách làm này đơn giản và luôn
nhất quán với trạng thái hiện tại, còn chi phí của nó được đo ở Mục~7.

\paragraph{Lật lời khai}
Khi một phán quyết được chấp nhận, hàm \texttt{\_apply\_proof(state, i, s)} hiện
thực quy tắc \eqref{eq:reveal} qua hai bước: ghi \texttt{state.proved\_statuses[i]
= s}, sau đó thêm vào \texttt{state.revealed\_clue\_ids} định danh của mọi lời
khai có \texttt{source\_cell\_id} bằng $i$. Bước lật không phụ thuộc vào giá trị
$s$, do đó lời khai của tội phạm cũng được lật như lời khai của người vô tội. Vì
ảnh chụp trả về sau \texttt{ACCEPTED} đã chứa lời khai mới, lời gọi kế tiếp, dù
đến từ người chơi hay tác nhân, đều có thể sử dụng ngay thông tin này. Ca kiểm
thử \texttt{test\_revealed\_clue\_makes\_new\_cell\_provable} kiểm tra đúng chuỗi
sự kiện đó, khi một ô ban đầu nhận \texttt{NOT\_PROVABLE} trở thành
\texttt{ACCEPTED} sau khi lời khai mới được lật.

\paragraph{Hoàn tất ván}
Engine không có trạng thái ``kết thúc'' riêng; ván được xem là hoàn tất khi
\texttt{PublicState.is\_complete} đúng, tức là mọi ô đều đã có mặt trong
\texttt{proved\_statuses}. Kể từ thời điểm đó, mọi lần nộp tiếp theo đều nhận
\texttt{ALREADY\_PROVED} nên trạng thái ván không thể thay đổi thêm. API trả về
trường \texttt{is\_complete} trong mỗi phản hồi để giao diện hiển thị màn hình
kết thúc, còn \texttt{DeductiveAgent.next\_move} trả về \texttt{SOLVED} khi không
còn ô nào chưa được chứng minh.

\paragraph{Khởi động lại}
Hàm \texttt{restart()} thay \texttt{self.\_state} bằng một kết quả mới của
\texttt{\_initial\_state()} và trả về ảnh chụp tương ứng. Do trạng thái được dựng
lại từ \texttt{PuzzleDefinition} bất biến, ván chơi luôn trở về đúng
$(\mathcal{R}_0, \mathcal{P}_0)$ bất kể trước đó đã đi bao nhiêu nước, trong khi
các ảnh chụp cũ đã phát ra vẫn không bị ảnh hưởng vì mỗi ảnh chụp giữ bản sao
riêng. Ở phía API, thao tác nạp bản đồ khác (\nolinkurl{/api/game/load}) không
gọi \texttt{restart} mà tạo hẳn một \texttt{GameEngine} mới rồi thay vào
\texttt{app.state.engine}.
\end{sloppypar}
